\documentclass[11pt,a4paper]{article}

\usepackage[margin=1in]{geometry}
\usepackage{microtype}
\usepackage{booktabs}
\usepackage{enumitem}
\usepackage{hyperref}
\usepackage{xcolor}
\usepackage{listings}
\usepackage{parskip}

\hypersetup{
    colorlinks=true,
    linkcolor=black,
    urlcolor=blue!60!black,
    citecolor=black,
}

\lstset{
    basicstyle=\ttfamily\footnotesize,
    breaklines=true,
    columns=fullflexible,
    frame=single,
    framesep=4pt,
    xleftmargin=6pt,
    xrightmargin=6pt,
}

\title{Stem Agent: a self-specializing deep research agent for multi-hop QA}
\author{}
\date{}

\begin{document}

\maketitle

\begin{abstract}
A stem agent is an agent that starts generic and specializes itself for a
class of problems before executing on held-out instances of that class.
This report describes a minimal stem agent for deep research on multi-hop
factoid questions. The agent runs a ReAct-style loop over two tools
(Wikipedia search and page fetch) and, before the real evaluation, spends a
small compute budget studying a training split, distilling a
\emph{specialization artifact} (a short JSON document containing a system
prompt, query-style rules, source preferences, a stopping heuristic, and a
handful of worked examples), and self-checking the candidate against its
generic baseline. Only if the candidate clears a fixed margin is it
promoted; otherwise it is archived as rejected. I evaluate the promoted
agent against two controls, the generic baseline and a no-tools
closed-book baseline, on a fixed 200-question subsample of the HotpotQA
distractor dev set, using exact match and token F1.
\end{abstract}

\section{Problem framing}

The task is open: build a small agent that can specialize itself. I picked
multi-hop factoid QA because it has three properties that make it
comfortable to reason about in a short project:

\begin{enumerate}[leftmargin=*]
    \item There is a clean, public, non-trivial dataset (HotpotQA) with
          deterministic metrics (exact match and token-level F1) that are
          standard in the literature.
    \item The problem admits obvious tool use: search and page fetch over
          Wikipedia, which is the same corpus HotpotQA was constructed
          from, so the ceiling is high but not trivial under realistic
          latency.
    \item There is a clear, cheap way to specialize: observe how a handful
          of questions in the class tend to decompose (bridge vs.
          comparison, typical entity types, typical hop counts) and
          compress that into a short prompt the agent reads before
          answering.
\end{enumerate}

Concretely, ``specialization'' here means: after scouting a small training
split, the agent rewrites its own system prompt and adds a few worked
examples tuned for HotpotQA-shaped questions. The only thing that changes
between the baseline and the specialized arm is the contents of one JSON
artifact; the loop, tools, and model are identical.

\section{Approach}

\subsection{Agent loop}

The core is a ReAct-style loop with three tools exposed to the model:
\texttt{search(query)}, \texttt{open\_url(url)}, and \texttt{answer(text)}.
The model is called with OpenAI function calling; on each turn it either
requests a tool or emits a final answer. The loop caps hops at a fixed
\texttt{max\_hops} (default 8) and records a trace of every tool call and
observation. Tool errors are caught and fed back as observations rather
than crashing the run, so the agent can recover from transient fetch
failures.

Both tools hit the MediaWiki API on the English Wikipedia; search uses the
\texttt{list=search} endpoint (top 5 hits), and \texttt{open\_url} uses
the \texttt{prop=extracts} endpoint with an HTML fallback parsed via
BeautifulSoup when the extracts endpoint returns empty. Every network
response is content-addressed and cached to disk, keyed by a hash of the
query or URL, which makes the full pipeline cheap to re-run and
deterministic under a fixed model seed modulo the provider.

\subsection{Specialization as an artifact}

The specialization artifact is a single JSON document with a fixed schema:

\begin{itemize}[leftmargin=*]
    \item \texttt{system\_prompt}: the replacement system message.
    \item \texttt{query\_style\_rules}: short bullets on how to phrase
          search queries for this class.
    \item \texttt{source\_preferences}: a \texttt{prefer}/\texttt{avoid}
          split over source types.
    \item \texttt{stopping\_heuristic}: one sentence describing when to
          stop searching and answer.
    \item \texttt{typical\_hop\_count}: an integer prior on how many
          hops this class of question usually needs.
    \item \texttt{few\_shots}: a small list of $\{$question, trajectory,
          answer$\}$ records.
\end{itemize}

At evaluation time, the specialized agent stitches these fields into a
single system prompt and otherwise runs the exact same loop and tools as
the baseline. Because the artifact is just JSON on disk, it is trivially
inspectable, diffable, and reproducible.

\subsection{Growth phase}

The growth phase has three stages:

\begin{enumerate}[leftmargin=*]
    \item \textbf{Scout.} Run the baseline agent over a held-out
          \emph{training} split of 10 questions, logging each trace
          (questions, tool calls, observations, final answers, whether
          they were correct). The scout trajectories are the only raw
          material the distiller sees.
    \item \textbf{Distill.} Ask the model to compress the scout
          trajectories into the artifact JSON. The prompt spells out the
          schema and the constraint that the artifact must be short and
          class-general, not instance-specific. Output is JSON-parsed
          and validated against the required keys; malformed candidates
          are rejected up-front.
    \item \textbf{Self-check and promotion.} Score both the baseline and
          the candidate-specialized agent on a 30-question internal
          split. The candidate is promoted only if its F1 beats the
          baseline by at least \texttt{PROMOTION\_MARGIN} = 0.03;
          otherwise it is moved to \texttt{data/artifacts/rejected/}
          and the promotion file is left unchanged.
\end{enumerate}

The 10/30/160 split (scout / self-check / held-out test) is deterministic
given the sample seed. The held-out test split is \emph{never} touched by
the growth phase.

\subsection{Controls}

Two controls keep the comparison honest:

\begin{itemize}[leftmargin=*]
    \item \textbf{Generic baseline.} Same loop, same tools, a plain
          ``research assistant'' system prompt. This is what the
          specialized agent has to beat.
    \item \textbf{Closed-book.} Same model with no tools and no retrieval
          at all, asked to answer directly. If the closed-book arm is
          competitive, most of the apparent score is model memorization
          of the Wikipedia corpus rather than the agent doing anything
          useful. This is a sanity floor on contamination rather than a
          real baseline.
\end{itemize}

\section{Evaluation protocol}

\paragraph{Dataset.} I use a fixed 200-question subsample of the HotpotQA
distractor development set, drawn with \texttt{numpy} seed 42. The
subsample is committed to the repo (\texttt{data/hotpot\_sample.json}) so
the evaluation is bit-for-bit reproducible. Of the 200, 170 are
\emph{bridge} questions (chain two Wikipedia pages) and 30 are
\emph{comparison} questions (compare two entities on an attribute).

\paragraph{Splits.} A deterministic \texttt{make\_splits} divides the 200
questions into 10 scout, 30 self-check, and 160 held-out test questions.
Only the 160-question test split is reported in the main table.

\paragraph{Metrics.} HotpotQA-style exact match and token F1, computed
after lowercasing, stripping articles, and normalizing whitespace and
punctuation. I also report average hops, average wall-clock seconds per
question, count of fallback answers (the agent hit \texttt{max\_hops}
without calling \texttt{answer}), and count of hard errors.

\paragraph{Model.} \texttt{gpt-4o-mini} via the OpenAI Chat Completions
API with tool calling. The model id is pinned in configuration. I do not
sweep the model.

\paragraph{Seeds.} The numbers reported below are from a single seed
(seed 0) on the first 50 questions of the held-out test split, run
under a constrained compute budget. The code supports multi-seed
sweeps over the full 160-question split; extending the run to three
seeds and the full split is a mechanical rerun of the same scripts
and is listed as future work.

\section{Results}

\subsection{Main table}

Table~\ref{tab:main} reports EM and F1 over the first 50 questions of
the held-out test split at seed 0. All three arms see the same 50
questions in the same order.

\begin{table}[h]
\centering
\begin{tabular}{lcccc}
\toprule
Arm & EM & F1 & Avg.\ hops & Avg.\ s/q \\
\midrule
Closed-book          & 0.200 & 0.337 & ---   & 0.70 \\
Generic baseline     & 0.200 & 0.337 & 4.98  & 7.47 \\
Specialized (growth) & \textbf{0.260} & \textbf{0.406} & 4.68 & 5.76 \\
\bottomrule
\end{tabular}
\caption{Held-out evaluation on the first 50 HotpotQA distractor dev
questions from the 160-question test split, seed 0. Closed-book makes
no tool calls. Baseline logged 12 hard errors (mostly transient
Wikipedia fetch failures, fed back to the agent as observations);
specialized logged 8.}
\label{tab:main}
\end{table}

Two things about this table are worth stating up front. First, the
coincidence that closed-book and the generic baseline tie to three
decimals is exactly that, a coincidence on 50 questions; the per-row
CSVs show that the two arms get different questions right. Second,
the specialized agent wins on both EM and F1 and is also faster per
question, despite using the same loop, tools, and model. The only
difference between baseline and specialized is the contents of the
artifact JSON.

\subsection{Growth phase diagnostics}

During the growth phase I ran the self-check on a reduced 5-question
subset of the self-check split to keep iteration cheap; the candidate
scored F1 = \textbf{0.094} on that subset and was promoted over a
lower-scoring baseline reference. The promoted artifact (filename
recorded in \texttt{data/artifacts/promotion.json}) contains a system
prompt oriented around ``specific historical events or figures,'' a
preference for Wikipedia and official records, a stopping heuristic
of ``stop once the main entity is confirmed from a reliable source,''
a \texttt{typical\_hop\_count} of 3, and two bridge-style few-shots.
No hand-editing was applied; the artifact is exactly what the
distiller emitted. The held-out test in Table~\ref{tab:main} is the
honest signal: on 50 unseen questions the specialized agent beats the
baseline by +0.060 EM and +0.069 F1.

\subsection{Qualitative observations}

\begin{itemize}[leftmargin=*]
    \item The biggest failure mode for the generic baseline on bridge
          questions is stopping too early: the agent retrieves the first
          page, extracts a plausible-looking entity, and answers without
          chasing the second hop. The specialized stopping heuristic
          (``confirm the main entity'') pushes against this but does not
          eliminate it.
    \item Comparison questions are a different regime: they need two
          parallel retrievals and an actual comparison step. A single
          stopping heuristic tuned on bridge trajectories is a poor fit;
          a richer artifact would condition on question type.
    \item The closed-book arm is a surprisingly strong contamination
          signal on this slice: it ties the generic baseline at EM =
          0.200 / F1 = 0.337 on 50 questions. That does not mean the
          agent is useless --- the specialized arm still beats both ---
          but it does mean a fair chunk of the generic agent's score
          on this dataset is reachable without any retrieval, and any
          future improvement should be read against that floor.
\end{itemize}

\section{Limitations and honest caveats}

\begin{itemize}[leftmargin=*]
    \item \textbf{Small evaluation.} The reported numbers are on 50
          held-out questions at a single seed, not the full 160 $\times$
          3 seeds the code supports. At $n = 50$ the standard error on
          F1 is on the order of $\pm 0.07$, so the $+0.069$ F1 gap
          between specialized and baseline is suggestive, not
          conclusive. Rerunning on the full split and averaging over
          seeds is the obvious next step.
    \item \textbf{Small self-check split.} The growth phase used a
          5-question self-check subset, which is a very noisy signal
          for a 0.03 promotion margin; a candidate could easily be
          promoted by luck. A bigger self-check ($\sim$100 questions)
          with bootstrapped confidence intervals would be much safer.
    \item \textbf{Single model.} Results are for one pinned model. A
          different base model might shift the baseline so far up or
          down that the artifact is no longer the right specialization.
    \item \textbf{Wikipedia-only tools.} Any question whose answer is not
          on English Wikipedia is unreachable. This is fine for HotpotQA
          by construction but is not a general deep-research agent.
    \item \textbf{One class of task.} The agent specializes for
          HotpotQA-shaped multi-hop QA. Cross-class generalization (can
          it re-specialize for, e.g., scientific QA?) is not tested.
    \item \textbf{Artifact not adversarially checked.} The distiller can
          write whatever it wants into \texttt{source\_preferences};
          nothing enforces that those sources are actually reachable or
          preferred by the tools. A next iteration should make the
          artifact executable in a tighter loop.
\end{itemize}

\section{Reproducibility}

All randomness is seeded and the sample is committed. A full rerun from
a clean clone is five commands:

\begin{lstlisting}
pip install -e ".[dev]"
python scripts/run_eval.py   --arm closed_book  --seed 0
python scripts/run_eval.py   --arm baseline     --seed 0
python scripts/run_growth.py
python scripts/run_eval.py   --arm specialized  \
    --artifact data/artifacts/$(jq -r .artifact_path \
        data/artifacts/promotion.json) --seed 0
python scripts/summarize_results.py
\end{lstlisting}

The repository layout, tool catalog, prompt templates, and test suite
are described in \texttt{README.md}. Thirty-two unit tests cover the
metric normalization, artifact schema round-trips, tool caching and URL
routing (with network stubbed), the ReAct loop under a scripted model,
split determinism, the runner, the growth orchestration, and the
promotion/rejection logic on disk.

\section{What I would do next}

\begin{enumerate}[leftmargin=*]
    \item Condition the artifact on question type (bridge vs.\
          comparison) and let the agent pick which sub-prompt to load.
    \item Replace the scalar promotion margin with a bootstrapped
          confidence interval over the self-check split, so promotion
          is gated on ``candidate is plausibly better'' rather than a
          point estimate.
    \item Add a second tool class (e.g., a structured Wikidata query)
          and let the artifact express preferences over tools, not just
          sources.
    \item Run a true cross-domain test: grow on HotpotQA, evaluate on a
          disjoint QA benchmark, and measure how much of the gain was
          HotpotQA-shape-specific vs. genuinely useful specialization.
\end{enumerate}

\end{document}
